\documentclass[11pt,a4paper]{article}
\usepackage[margin=18mm]{geometry}
\usepackage{fontspec}
\setmainfont{DejaVu Serif}
\setsansfont{DejaVu Sans}
\setmonofont{DejaVu Sans Mono}
\usepackage{polyglossia}
\setdefaultlanguage{russian}
\usepackage{array,longtable,booktabs,tabularx}
\usepackage{xcolor}
\usepackage{enumitem}
\usepackage{hyperref}
\usepackage{fancyhdr}
\usepackage{titlesec}
\usepackage{listings}
\usepackage{microtype}
\usepackage{parskip}
\definecolor{AltairBlue}{HTML}{0E5A88}
\definecolor{AltairGray}{HTML}{4B5563}
\hypersetup{colorlinks=true,linkcolor=AltairBlue,urlcolor=AltairBlue,pdftitle={ЦУП Альтаир v7.1 — руководство пользователя}}
\pagestyle{fancy}
\fancyhf{}
\lhead{ЦУП Альтаир v7.1}
\rhead{Технопром 2026 · Миссия на Луну}
\cfoot{\thepage}
\titleformat{\section}{\Large\bfseries\sffamily\color{AltairBlue}}{\thesection.}{0.5em}{}
\titleformat{\subsection}{\large\bfseries\sffamily}{\thesubsection.}{0.5em}{}
\lstset{basicstyle=\ttfamily\small,breaklines=true,frame=single,columns=fullflexible,showstringspaces=false}
\newcommand{\code}[1]{\texttt{\detokenize{#1}}}

\begin{document}

\begin{center}
{\Huge\bfseries\sffamily ЦУП Альтаир v7.1}\\[2mm]
{\large Центр управления полётами}\\[3mm]
{\Large\bfseries Технопром 2026 · Миссия на Луну}\\[8mm]
{\large Руководство по установке, прошивке и работе с телеметрией}\\[8mm]
\hrule
\end{center}

\vspace{6mm}

\textbf{Назначение.} ЦУП Альтаир — локальное Windows-приложение учебной наземной станции. В рабочей конфигурации STM32 формирует и передаёт пакет телеметрии через CC1101, ESP32 с CC1101 принимает пакет на частоте 435 МГц и передаёт данные по USB в ЦУП.

\textbf{Проверенная конфигурация v7.1.} На стенде устойчивый приём получен при широкой полосе приёмника \textbf{203 кГц}. В версии 0.7.1 встроенная прошивка ESP32 приведена к тому же алгоритму, который реально заработал при ручной загрузке через Arduino IDE: \code{radio.receive(packet)} без старой IRQ-обёртки и без короткого пользовательского таймаута 250 мс.

\textbf{Исправление относительно 0.7.0.} В релизе 0.7.0 ручной файл \code{Altair_Gateway_v7.ino} и встроенный \code{gateway-firmware.bin} собирались из разных реализаций приёмника. Поэтому ручной скетч работал, а прошивка из установщика могла не выдавать телеметрию. В 0.7.1 исходники синхронизированы, а CI блокирует выпуск, если они снова отличаются.

\section{Что скачать}

На главной странице репозитория OpenMCC размещены прямые ссылки на четыре основных файла:

\begin{enumerate}[leftmargin=7mm]
\item \code{CUP-Altair-Setup-0.7.1.exe} — программа ЦУП Альтаир для Windows.
\item \code{Altair_Gateway_v7.ino} — скетч для ESP32-WROOM-32 + CC1101.
\item \code{Transmit_v7_static.ino} — тестовый скетч передатчика для STM32 + CC1101.
\item \code{CUP_Altair_v7_Manual.pdf} — настоящее руководство.
\end{enumerate}

\section{Архитектура системы}

\begin{center}
\fbox{\parbox{0.92\linewidth}{\centering
\textbf{STM32 + CC1101} \quad $\longrightarrow$ \quad \textbf{435.000 МГц} \quad $\longrightarrow$ \quad \textbf{CC1101 + ESP32-WROOM-32} \quad $\longrightarrow$ \quad \textbf{USB 115200} \quad $\longrightarrow$ \quad \textbf{ЦУП Альтаир v7.1}
}}
\end{center}

Передатчик и приёмник используют одинаковые параметры пакетного режима CC1101. Отличие полосы приёма не мешает передаче: широкая полоса 203 кГц применяется на наземном приёмнике для устойчивой работы при заметном частотном уходе модулей.

\section{Подключение CC1101 к ESP32}

\begin{center}
\begin{tabular}{lll}
\toprule
\textbf{CC1101} & \textbf{ESP32-WROOM-32} & \textbf{Назначение}\\
\midrule
VCC & 3V3 & питание 3.3 В\\
GND & GND & общий провод\\
SCK & GPIO18 & SPI clock\\
MISO / SO & GPIO19 & SPI MISO\\
MOSI / SI & GPIO23 & SPI MOSI\\
CSN / CS & GPIO21 & chip select\\
GDO0 & GPIO4 & окончание принятого пакета\\
GDO2 & GPIO27 & служебный вывод CC1101\\
\bottomrule
\end{tabular}
\end{center}

\textbf{Важно:} питание CC1101 — только 3.3 В. Перед включением необходимо проверить общий GND и правильность SPI-подключения.

\section{Подключение CC1101 к STM32}

В тестовом передатчике сохранено подключение, с которым ранее была получена работа радиомодуля IntroSat:

\begin{center}
\begin{tabular}{ll}
\toprule
\textbf{Сигнал CC1101} & \textbf{STM32}\\
\midrule
CS / CSN & PA2\\
GDO0 & PA0\\
GDO2 & PA1\\
SCK / MISO / MOSI & штатные SPI-выводы выбранной STM32-платы\\
\bottomrule
\end{tabular}
\end{center}

Если используется другая STM32-плата, следует сохранить фактическое рабочее подключение и выбрать соответствующую плату в Arduino IDE / PlatformIO.

\section{Радиопрофиль v7.1}

\begin{center}
\begin{tabular}{ll}
\toprule
Параметр & Значение\\
\midrule
Частота & 435.000 МГц\\
Модуляция & 2-FSK\\
Скорость & 4.8 kbps\\
Девиация & 5 кГц\\
Полоса RX наземного приёмника & \textbf{203 кГц}\\
Sync word & 0x12AD\\
Кодирование & NRZ\\
Длина пакета & variable\\
CRC CC1101 & включён\\
Преамбула & 16 bit\\
Мощность тестового передатчика & 5 dBm\\
\bottomrule
\end{tabular}
\end{center}

Полоса 203 кГц выбрана не как теоретический минимум, а как \textbf{рабочая стендовая настройка}: при 58 кГц пакет не принимался, а при 203 кГц тот же пакет уверенно декодировался.

\section{Формат пакета телеметрии}

Статический тестовый пакет:

\begin{lstlisting}
02,00001,00015,3.00,4.20,1,33
\end{lstlisting}

Фактическая длина — \textbf{29 ASCII-символов}. Порядок полей:

\begin{center}
\begin{tabularx}{\linewidth}{>{\bfseries}l c X}
\toprule
Поле & Длина & Смысл\\
\midrule
ID & 2 & номер команды / идентификатор аппарата\\
PACKET & 5 & порядковый номер пакета\\
UPTIME & 5 & время с момента запуска, с\\
PANEL\_POWER & 4 & мощность солнечных панелей, X.XX Вт\\
VOLT & 4 & напряжение аккумулятора, X.XX В\\
MODE & 1 & 1 — штатный, 0 — аварийный\\
CHECKSUM & 2 & XOR в HEX\\
\bottomrule
\end{tabularx}
\end{center}

Аварийный режим в окончательной бортовой программе должен включаться автоматически при \code{VOLT < 2.00 В}.

\subsection{Контрольная сумма XOR}

XOR рассчитывается по всей строке до двух HEX-символов контрольной суммы, \textbf{включая последнюю запятую}:

\begin{lstlisting}
02,00001,00015,3.00,4.20,1,
\end{lstlisting}

Для приведённого примера результат равен \code{0x33}.

\section{Что показывает ЦУП}

Карточки телеметрии расположены в следующем порядке:

\begin{enumerate}[leftmargin=7mm]
\item ID спутника;
\item номер пакета;
\item время работы с момента запуска;
\item мощность солнечных панелей;
\item напряжение аккумулятора;
\item режим работы;
\item контрольная сумма XOR;
\item раскрытие антенны;
\item RSSI;
\item SNR.
\end{enumerate}

В минимальном 29-символьном пакете состояния антенны пока нет. Поэтому ЦУП v7.1 показывает в этой карточке \textbf{«Н/Д»}, а не прочерк. Когда в будущий формат будет добавлено поле антенны, значения должны трактоваться как \code{1 = РАСКРЫТА}, \code{0 = СЛОЖЕНА}.

RSSI и SNR всегда располагаются \textbf{в самом конце, после карточки раскрытия антенны}. RSSI берётся из наземного CC1101. SNR для CC1101 является диагностической оценкой относительно принятого уровня шумового фона / опорного уровня и не является отдельным аппаратным измерением CC1101.

\section{Графики}

В v7.1 строятся графики:

\begin{itemize}[leftmargin=7mm]
\item напряжение аккумулятора;
\item мощность солнечных панелей;
\item RSSI;
\item SNR.
\end{itemize}

Каждый отсчёт отображается \textbf{точкой на линии}. Это позволяет видеть фактические моменты поступления пакетов, а не только сглаженную кривую.

\section{Прошивка STM32}

\begin{enumerate}[leftmargin=7mm]
\item Открыть \code{Transmit_v7_static.ino} в Arduino IDE.
\item Установить библиотеку RadioLib.
\item Выбрать фактическую STM32-плату и способ загрузки.
\item Проверить рабочие выводы CC1101: PA2 / PA0 / PA1 и штатный SPI.
\item Загрузить скетч.
\item В Serial Monitor убедиться, что раз в секунду появляется \code{TX OK}.
\end{enumerate}

Тестовый скетч передаёт статическую строку \code{02,00001,00015,3.00,4.20,1,33} раз в секунду.

\section{Прошивка ESP32}

\subsection{Вариант A — через Arduino IDE}

Это наиболее прозрачный способ для первичной проверки.

\begin{enumerate}[leftmargin=7mm]
\item Открыть \code{Altair_Gateway_v7.ino}.
\item Выбрать ESP32-WROOM-32 / совместимую плату.
\item Установить RadioLib.
\item Загрузить скетч.
\item Открыть Serial Monitor, 115200 бод.
\item При работающем STM32 должна появиться строка вида:
\end{enumerate}

\begin{lstlisting}
$TEL,ID=02,PACKET=00001,UPTIME=00015,PANEL_POWER=3.00,VOLT=4.20,MODE=1,CHECKSUM=33,CHECKSUM_OK=1,RSSI=-23.5,SNR=60.0,LQI=52
\end{lstlisting}

После проверки \textbf{закрыть Serial Monitor}, иначе он будет удерживать COM-порт.

\subsection{Вариант B — встроенная прошивка из ЦУП Альтаир}

Установщик v7.1 содержит скомпилированную прошивку радиошлюза. В версии 0.7.1 её приёмный тракт синхронизирован с \code{Altair_Gateway_v7.ino}: используется тот же вызов \code{radio.receive(packet)} и та же полоса 203 кГц. Поэтому ручная прошивка и прошивка из программы должны давать одинаковое поведение радиоприёмника.

Если ESP32 уже прошита через Arduino IDE рабочим скетчем, для проверки телеметрии повторная прошивка из ЦУПа не требуется. Если используется встроенный прошивальщик v7.1, после его завершения ESP32 следует перезапустить и затем подключить основной COM-порт.

\section{Запуск ЦУП Альтаир}

\begin{enumerate}[leftmargin=7mm]
\item Установить \code{CUP-Altair-Setup-0.7.1.exe}.
\item Подключить прошитую ESP32 к компьютеру по USB.
\item Закрыть Arduino Serial Monitor и другие программы, удерживающие COM-порт.
\item Запустить «ЦУП Альтаир».
\item В разделе «Соединение» выбрать профиль ESP32 и скорость 115200 бод.
\item Нажать «Подключить устройство» и выбрать COM-порт ESP32.
\item Включить STM32-передатчик.
\item Убедиться, что карточки телеметрии обновляются, а на графиках появляются точки.
\end{enumerate}

\section{Быстрая диагностика}

\begin{longtable}{p{0.30\linewidth}p{0.62\linewidth}}
\toprule
\textbf{Симптом} & \textbf{Что проверить}\\
\midrule
\endfirsthead
\toprule
\textbf{Симптом} & \textbf{Что проверить}\\
\midrule
\endhead
ESP32 не видит CC1101 & питание 3.3 В, GND, GPIO18/19/23/21; простой тест \code{radio.begin()} должен вернуть 0.\\
Пакеты видны на SDR, но нет приёма & убедиться, что RX bandwidth = 203 кГц; проверить 435.000 МГц, 4.8 kbps, 2-FSK, deviation 5 кГц, sync 0x12AD, CRC.\\
Ручной скетч принимает, а прошивка из программы нет & использовать версию не ниже 0.7.1; в 0.7.0 встроенный binary и ручной скетч имели разные реализации RX.\\
В Serial Monitor есть \code{$TEL,...}, а ЦУП пустой & закрыть Serial Monitor и открыть тот же COM-порт в ЦУПе на 115200 бод.\\
ЦУП не открывает COM-порт & закрыть Arduino IDE Serial Monitor, SDR-сервис с Serial, терминалы и другие программы, которые используют порт.\\
Антенна показывает «Н/Д» & это штатно для 29-символьного пакета v7: поля состояния антенны в нём пока нет.\\
Сигнал на SDR смещён от 435.000 МГц & широкая полоса 203 кГц введена именно для повышения устойчивости к такому рассогласованию; при сильном смещении дополнительно проверить кварцы и фактическую частоту обоих модулей.\\
\bottomrule
\end{longtable}

\section{Минимальный контроль перед демонстрацией}

Перед работой на площадке рекомендуется выполнить полный цикл:

\begin{enumerate}[leftmargin=7mm]
\item STM32 передаёт пакет раз в секунду — подтверждено Serial Monitor / SDR.
\item ESP32 принимает пакет с BW 203 кГц и выдаёт \code{$TEL,...}.
\item XOR пакета равен \code{33} и отображается как OK.
\item ЦУП получает ID, пакет, uptime, мощность, напряжение и режим.
\item Карточка антенны показывает «Н/Д», RSSI и SNR находятся после неё.
\item На четырёх графиках появляются линии с точками.
\end{enumerate}

\vfill
\begin{center}
\small ЦУП Альтаир v7.1 · Технопром 2026 · Миссия на Луну
\end{center}

\end{document}
